\section{可換 AOU 点の形式変形と集合論的な標識変更}
\label{sec:formal-abelian-recognition}
本節は固定された可換反標準 heart の変形理論の補助結果を記す。
得られる形式 Abel 変形が反標準 heart の許容条件を満たす範囲を、
抽象的な Abel 変形の全体と区別する。
鍵になるのはDe Deken--Lowen\cite[Theorems 4.5, 5.13, 5.15]{DDL}であり、
本稿の有限代数変形の完備性（定理\ref{thm:aou-relation-formal-completeness}）
と組み合わせる。導来圏の任意のheartを自動的に認識するのではなく、
特定の可換三次曲面から始まる\emph{Artin形式変形}を扱う。

\subsection{標準helixとAbel変形の比較}
$X$を滑らかな複素三次曲面とし、
$E_0,\ldots,E_8$をAOUの標準九対象line-bundle threadとする。
$\mathscr H=(E_i)_{i\in\mathbb Z}$は、このthreadが生成する
幾何学的helixである。以下、文献でいう
\emph{$\mathbb Z$-generating geometric helix}の条件、すなわち
$(E_i)$が$\operatorname{Qcoh}X$を$\mathbb Z$-生成し、
$i\le j$について$\operatorname{Ext}^{1,2}(E_i,E_j)=0$を満たすことを用いる。
標準三ブロックhelixの幾何学的性質は
Karpov--Nogin\cite{KN}およびAOU\cite[\S4]{AOU}の既知の入力である。
$-K_X$の豊富性と$E_{i+9}\cong E_i\otimes\mathcal O_X(-K_X)$により、
このhelixの適切な反標準方向の部分列が生成列となる。

\begin{theorem}[AOU形式関係式からのAbel変形の再構成]
\label{thm:aou-formal-abelian-reconstruction}
$X,\mathscr H$を上のとおりとし、$r\in R_Q$をそのAOU関係式点、
$B=\operatorname{End}_X(\bigoplus_{i=0}^{8}E_i)$とする。
$R$を剰余体$\mathbb C$の局所Artin $\mathbb C$-代数とする。
$r$に特殊化する任意のAOU関係式の$R$-点$r_R$について、
付随する平坦な有限代数$B_R$から、
$\mathcal C_R\otimes_R\mathbb C\simeq\operatorname{Qcoh}X$という
特殊ファイバー同定を伴う平坦なGrothendieck Abel変形
$\mathcal C_R$が構成できる。
標準threadの変形を保持したAbel変形の同型類集合と、
九つの原始冪等元を標識した$B$の平坦な有限代数変形の同型類集合は
自然に全単射である。これらの同値は$R$の小拡大と底変更に両立する。
さらに$\mathcal C_R$は局所noetherianであり、
そのnoetherian対象の部分圏は$\operatorname{coh}X$からの
Abel変形を与える。
\end{theorem}
\begin{proof}
標準helixから作る$\mathbb Z$-代数を$\mathfrak a$、
九対象threadへの制限を$\mathfrak a_{[0,8]}$と書く。
後者は$B$の原始冪等元を標識した線形カテゴリーである。
De Deken--Lowen\cite[Theorem 5.15]{DDL}は、上で述べた生成性、
幾何学的helixおよびExt消滅のもとで
\[
 \operatorname{Def}_{\rm ab}(\operatorname{Qcoh}X)
 \simeq\operatorname{Def}_{\rm lin}(\mathfrak a)
 \simeq\operatorname{Def}_{\rm lin}(\mathfrak a_{[0,8]})
\]
を与える。第二の変形問題は、原始冪等元を標識した
有限代数$B$の平坦な変形問題と等しい。
定理\ref{thm:aou-relation-formal-completeness}により
この変形は九つの関係式の平坦な変形により尽くされる。
上記の同値の逆を適用すれば$\mathcal C_R$を得る。
この議論はAbel変形の同値を使うため、
単なる$D^b(\operatorname{mod}B_R)$の標準heartを
$\operatorname{coh}X$と取り違えてはいない。

局所noetherian性も確認する。De Deken--Lowenの再構成では
$\mathbb Z$-生成列$(E_i)$が$R$-flatな生成対象$(E_{i,R})$へ持ち上がる。
$R$の極大イデアルを$\mathfrak m$とすると
$E_{i,R}$の有限$\mathfrak m$-adic filtrationの各商は
$E_i$の有限直和である。$E_i$はnoetherianであるから
$E_{i,R}$もnoetherianである。
従って$\mathcal C_R$はnoetherian生成対象を持つ
局所noetherian Grothendieck圏であり、noetherian対象の部分圏は
核・余核・有限拡大で閉じる。特殊ファイバーのnoetherian対象は
$\operatorname{coh}X$となる。
\end{proof}

\begin{remark}[復元の適用範囲]
\label{rem:aou-formal-heart-scope}
この定理は\emph{滑らかな可換点を固定した任意の局所Artin底}での
Abel変形の存在を与える。$R_Q$の全閉点で三次曲面型のheartが
代数的族として存在するという主張ではない。
またGrothendieck Abel変形の存在だけから
$\mathsf S_R^{-1}[2]$の$t$-exact性、rank-one構造対象の公理、
Artin--Zhang豊富性が形式変形の全域で自動的に成立すると
断言しない。これらを要する本稿の弱標識付きモジュライへの
因子化には、さらに検証を要する。
\end{remark}

\subsection{形式八次元チャートと標識変更}
$\mathcal D_X$を$\operatorname{Qcoh}X$の平坦なAbel変形関手で、
特殊ファイバーとの同定を保持し、
関手の自然同型を同一視したものとする。

\begin{theorem}[AOUの形式八次元チャート]
\label{thm:aou-formal-chart}
上の$X,r$について、$r$はAOUの全係数非零のトーラス
$M_Q^\circ$の点であるとする。
標準helixの選択に依存する形式変形関手の同型
\[
 \widehat{M_Q^\circ}_{[r]}\simeq \mathcal D_X
\]
が存在する。従って特殊ファイバーを同定したAbel変形の
同型類は、$\mathbb C$上の八次元形式円板で局所的に記述される。
その接空間は
\[
 T_{[r]}M_Q^\circ\simeq HH^2(B)\simeq HH^2(X)
 =H^1(X,T_X)\oplus H^0(X,-K_X)
\]
である。形式チャートの同型は標準helixと特殊化の同定に依存し、
無標識圏の大域的な正則パラメータ表示ではない。
\end{theorem}
\begin{proof}
定理\ref{thm:aou-formal-abelian-reconstruction}により
Abel変形は有限代数変形と同値であり、
定理\ref{thm:aou-relation-formal-completeness}により
九関係式の変形を矢の基底変更で割ったものと同値である。
全係数非零のトーラスでは、矢の基底変更群の
関係式への作用の核は頂点スカラーによる八次元トーラスである。
有効商トーラスの作用は自由である。
実際、作用行列のSmith標準形が
$\operatorname{diag}(1^{\times10},0^{\times8})$となることは、
命題\ref{prop:aou-kodaira-spencer-iso}の格子計算からも分かる。
従って関係式変形の同型類関手は
滑らかなトーラス商$M_Q^\circ$の点$[r]$における形式完成となる。
接空間については命題\ref{prop:aou-kodaira-spencer-iso}を適用する。
\end{proof}

\begin{proposition}[形式Picard作用と一次shear]
\label{prop:formal-picard-action}
$X$のexact自己同値群
$G_X=\operatorname{Pic}(X)\rtimes\operatorname{Aut}(X)$は、
特殊ファイバーの同定を変更することにより
$\mathcal D_X$に作用する。前定理の形式チャートへ輸送すると、
$\widehat{M_Q^\circ}_{[r]}$上に
一般には非線形な\emph{形式的}作用が定まる。
その$L\in\operatorname{Pic}(X)$による接空間作用は、
方向の規約による符号を除いて
\[
 (\xi,\pi)\longmapsto
 (\xi+H^1(\pi^\sharp)c_1(L),\pi)
\]
である。これは命題\ref{prop:picard-first-order}のshearに等しい。
この形式作用からAOUの大域GITスキームへのPicard群の
正則作用は導かれない。
\end{proposition}
\begin{proof}
特殊ファイバーとの同定を$\phi$から$F\circ\phi$へ変更する操作は
各$F\in\operatorname{Aut}_{\rm ex}(\operatorname{Qcoh}X)$について
Abel変形関手の自己同型を与え、合成に関して群作用となる。
可換滑らかな射影多様体のAbel圏の自己同値は
自己同型の引戻しと線束捩れで記述され、
$G_X=\operatorname{Pic}(X)\rtimes\operatorname{Aut}(X)$である。
定理\ref{thm:aou-formal-chart}によって作用を形式完成上に輸送する。
一次項は線束Morita変換のČech計算
（命題\ref{prop:picard-first-order}）により与えられる。
これはToda\cite{Toda}のHochschild変形に対する
Fourier--Mukai変換とも整合する。
\end{proof}

\begin{corollary}[無標識形式変形の同値類集合]
\label{cor:aou-formal-unmarked-quotient}
特殊ファイバーとの同定を忘れた Abel 変形の exact 同値類集合は、
同定付き変形の同型類集合 $\mathcal D_X$ の
$G_X=\operatorname{Pic}(X)\rtimes\operatorname{Aut}(X)$ による軌道集合である。
この主張は集合値の形式変形問題についてのものである。
反標準 heart と弱標識の条件を必要とする場合は、その条件を満たす部分に制限する。
\end{corollary}
\begin{proof}
基礎圏の二つの同定の差は、その特殊ファイバーの exact 自己同値である。
逆に同定の変更は、同定を忘れた圏の同値類を変えない。
従って同値関係はちょうどこの軌道関係である。
\end{proof}

\begin{remark}[共通線束捩れと形式チャートの非標準性]
\label{rem:formal-common-twist}
同じ$X$でthreadを$(E_i\otimes L)$へ取り替えても
$\operatorname{End}(\bigoplus_iE_i\otimes L)$は$B$と同一視され、
AOU関係式点は変わらない（命題\ref{prop:common-tensor}）。
しかしDe Deken--Lowenの逆構成を特殊ファイバー同定付きで比較すると、
二つの形式チャートは線束捩れによる再標識で移り合う。
その一次遷移は命題\ref{prop:formal-picard-action}のshearである。
従って関係式点から$\mathcal O$の絶対的な選択が自動的に
得られるわけではない。まさにこの差が無標識Abel圏を考える動機となる。
\end{remark}
